\documentclass[12pt,onecolumn]{IEEEtran}

\usepackage[utf8]{inputenc}
\usepackage[T1]{fontenc}
\usepackage[spanish]{babel}
\usepackage{amsmath}
\usepackage{graphicx}
\usepackage{hyperref}
\usepackage{cite}
\usepackage{geometry}

\geometry{margin=1in}

\title{Proyecto Final: Semáforo Vehicular y Peatonal}

\author{
Gabriel Alonso Chavarría Rodríguez\\
José Andrés Guerrero Araya\\
Julián Murillo Wo Ching\\
Jimena Vargas Vega
}

\date{\today}

\begin{document}

% The paper headers
\markboth{EL3215 Laboratorio de Electrónica Analógica, IS2026}%
{EL3215 Laboratorio de Electrónica Analógica}

\maketitle

\begin{abstract}
Resumen del proyecto.
\end{abstract}

\begin{IEEEkeywords}
Semáforo, FSM, temporización, lógica digital, electrónica analógica.
\end{IEEEkeywords}

%%%%%%%%%%%%%%%%%%%%%%%%%%%%%%%%%%%%%%%%%%%%%%%%%%%%%%%%%%%%%%
\section{Introducción}

Texto de introducción.

%%%%%%%%%%%%%%%%%%%%%%%%%%%%%%%%%%%%%%%%%%%%%%%%%%%%%%%%%%%%%%
\section{Propuesta de diseño}

\subsection{Subsistema de Control y Temporización}

\subsubsection{Reloj con Temporizador 555}

Texto.

\subsubsection{Botón y Antirrebote}

Texto.

\subsubsection{FSM}

Texto.

\subsubsection{Memoria del botón}

Texto.


%%%%%%%%%%%%%%%%%%%%%%%%%%%%%%%%%%%%%%%%%%%%%%%%%%%%%%%%%%%%%%
\section{Simulaciones}

\subsection{Simulación del reloj de 1 Hz}

Texto.

\subsection{Simulación del botón peatonal}

Texto.

\subsection{Simulación de la memoria de solicitud}

Texto.

\subsection{Simulación de salidas}

Texto.


%%%%%%%%%%%%%%%%%%%%%%%%%%%%%%%%%%%%%%%%%%%%%%%%%%%%%%%%%%%%%%
\section{Conclusiones}

\begin{itemize}
    \item Conclusión 1.
    \item Conclusión 2.
    \item Conclusión 3.
\end{itemize}

%%%%%%%%%%%%%%%%%%%%%%%%%%%%%%%%%%%%%%%%%%%%%%%%%%%%%%%%%%%%%%
\begin{thebibliography}{99}

\bibitem{fulano}


\end{thebibliography}

\end{document}